\section{The mean energy growth}
\label{sec:mean_energy_growth}
We finally introduce the mean energy growth which measures the amplification of perturbations averaged over an ensemble of random initial conditions. With the linearized operators (\S~\ref{sec:governing_equation}), the instantaneous disturbance energy (\S~\ref{sec:energy_norm}), and the statistics of the initial perturbation (\S~\ref{sec:statistics_of_the_incoming_perturbations}) now in place, all the ingredients required to define this quantity are available. At each wavenumber $k$, it is defined as the ratio of the expected energy at time $t$ to its initial value,
\begin{equation}
    \label{eq:mean_energy_growth_def}
    g(k,t)
    =
    \frac{\mean{\mathcal{E}_k(t)}}{\mean{\mathcal{E}_k(0)}}.
\end{equation}
The vertical part of the initial covariance enters as a discrete matrix \(\mathsfbi{C}_{00}^{z}\in\mathbb{C}^{2N_z\times 2N_z}\), the discretization of \(C_{00}^{z}\). Writing the perturbation state through the propagator, \(\hat{\mathbf{q}}_k(t)=\boldsymbol{\Phi}_{t,k}\hat{\mathbf{q}}_{0,k}\), and normalizing by \(\tr\{\mathsfbi{C}_{00}^{z}\}\), the purely vertical, wavenumber-resolved growth factor becomes
\begin{equation}
    g(k,t)
    =
    \frac{
      \tr
      \left\{
        \boldsymbol{\Phi}_{t,k}
        \mathsfbi{C}_{00}^{z}
        \boldsymbol{\Phi}_{t,k}^{*}
        \mathsfbi{W}_{k}
      \right\}
    }{
      \tr\{\mathsfbi{C}_{00}^{z}\}
    }.
    \label{eq:modal_energy_growth}
\end{equation}
The three-dimensional mean energy growth then follows by weighting this modal growth with the horizontal spectral density \(\Phi_{00}^{xy}\) and integrating over all wavenumbers,
\begin{equation}
    \Gm(t)
    =
    \frac{
      \displaystyle
      \iint
      g(k,t)
      \Phi_{00}^{xy}(k_x,k_y)
      \,\mathrm{d}k_x\,\mathrm{d}k_y
    }{
      \displaystyle
      \iint
      \Phi_{00}^{xy}(k_x,k_y)
      \,\mathrm{d}k_x\,\mathrm{d}k_y
    }.
    \label{eq:gmean_continuous_fourier}
\end{equation}

It is now clear that the finite-dimensional trace in \eqref{eq:modal_energy_growth} is the object that must be computed. The difficulty is that, for the non-autonomous semi-discrete system \eqref{eq:semidiscrete_system}, the propagator \(\boldsymbol{\Phi}_{t,k}\) is defined through time evolution and is not available as an explicitly formed matrix. If the system were autonomous, one could, at least in principle, compute the matrix exponential of the linear operator and evaluate the trace directly; although expensive, the calculation would then be straightforward. In the non-autonomous case considered here, the computation is more subtle, because the propagator is generated by the time-dependent matrices \(\mathsfbi{A}_{k}(t)\) and \(\mathsfbi{B}_{k}(t)\). The main algorithm developed in this work is therefore designed to compute this trace without ever explicitly forming the propagator matrix.


Because the linearized equation depends only on the magnitude \(k=\sqrt{k_x^2+k_y^2}\), the mean-energy formula \eqref{eq:gmean_continuous_fourier} can be written for this initialization in terms of the circularly integrated interface spectrum as
\begin{equation}
    \Gm(t)
    =
    \frac{
      \displaystyle
      \int_0^\infty
      g(k,t)\,
      E_{\eta}^{2D}(k)\,
      \mathrm{d}k
    }{
      \displaystyle
      \int_0^\infty
      E_{\eta}^{2D}(k)\,
      \mathrm{d}k
    } =
    \frac{2}{\mean{\eta_0^2}}
    \int_0^\infty
      g(k,t)\,
      E_{\eta}^{2D}(k)\,
      \mathrm{d}k.
    \label{eq:gmean_eta_spectrum}
\end{equation}
This expression clarifies the meaning of the three-dimensional energy growth. The per-mode mean energy \(g(k,t)\), defined in \eqref{eq:modal_energy_growth}, accounts for how the initial perturbations are correlated within the initial diffusive layer when a stochastic perturbation is imposed at the interface. For this class of problems, the resulting three-dimensional evolution is governed by the per-mode mean energy weighted by the spectrum that characterizes the interface disturbance. In other words, each wavenumber contribution is weighted by its corresponding contribution to the initial spectrum.

This observation implies that any growth rate based on classical modal analysis, or derived from the QSSA, is conditioned by the prescribed initial disturbances. The growth rate should therefore be defined for a specific initial spectrum and is expected to differ between different realizations. It is natural to define a mean growth rate for the full three-dimensional perturbation field as
\begin{equation}
    \label{eq:sigma_3d}
    \overline{\sigma}_{3D}(t)
    =
    \frac{1}{2}
    \difft{
        \log\!\left(\Gm(t)\right)
    }{t}.
\end{equation}
The corresponding one-dimensional growth rate associated with the per-mode mean energy growth is defined as
\begin{equation}
    \label{eq:sigma_1d}
    \overline{\sigma}_{1D}(k,t)
    =
    \frac{1}{2}
    \difft{
        \log\!\left(g(k,t)\right)
    }{t}.
\end{equation}

